\documentclass[11pt,a4paper]{article}
\usepackage[utf8]{inputenc}
\usepackage{amsmath,amssymb,amsfonts,amsthm}
\usepackage{bm}
\usepackage{geometry}
\geometry{margin=2.5cm}

\newtheorem{theorem}{Theorem}
\newtheorem{lemma}{Lemma}
\newtheorem{definition}{Definition}

\title{\textbf{Dual Equivalence Theorem: Whitney Edge Finite Elements and VNMM 2D $\mathcal{L}^1$}}
\author{Renato C. Mesquita}
\date{September 2026}

\begin{document}
\maketitle

\section{Introduction and Problem Setting}

Let $T \subset \mathbb{R}^2$ be a non-degenerate triangular simplex with vertices $V_1, V_2, V_3$ arranged in counter-clockwise order, and non-zero area $|T| > 0$.

\begin{definition}[Boundary Edges and Midpoints]
The three directed boundary edges of $T$ are denoted by $e_1 = (V_1, V_2)$, $e_2 = (V_2, V_3)$, and $e_3 = (V_3, V_1)$. For each edge $k \in \{1, 2, 3\}$:
\begin{itemize}
    \item Length: $L_k = \|V_{k+1} - V_k\| > 0$ (with cyclic index $V_4 \equiv V_1$);
    \item Unit tangent vector: $\mathbf{t}_k = \frac{V_{k+1} - V_k}{L_k}$;
    \item Edge midpoint: $M_k = \frac{1}{2}(V_k + V_{k+1})$.
\end{itemize}
\end{definition}

\begin{definition}[Whitney 1-Forms / N\'ed\'elec First Family]
Let $\lambda_i(\mathbf{x})$ be the affine barycentric coordinate associated with vertex $V_i$, such that $\lambda_i(V_j) = \delta_{ij}$ and $\sum_{i=1}^3 \lambda_i(\mathbf{x}) \equiv 1$. The standard Whitney 1-form associated with edge $e_k = (V_i, V_j)$ is:
\begin{equation}
    \mathbf{w}_k(\mathbf{x}) = \lambda_i(\mathbf{x}) \nabla \lambda_j - \lambda_j(\mathbf{x}) \nabla \lambda_i
\end{equation}
The polynomial approximation space is:
\begin{equation}
    \mathcal{W}(T) = \operatorname{span}\left\{ \begin{bmatrix} 1 \\ 0 \end{bmatrix}, \; \begin{bmatrix} 0 \\ 1 \end{bmatrix}, \; \begin{bmatrix} y \\ -x \end{bmatrix} \right\}
\end{equation}
The finite element degrees of freedom (DoFs) are edge circulations:
\begin{equation}
    \alpha_k(\mathbf{u}) = \oint_{e_k} \mathbf{u} \cdot d\boldsymbol{\ell} = \int_{e_k} (\mathbf{u} \cdot \mathbf{t}_k) \, ds, \quad \text{with} \quad \oint_{e_j} \mathbf{w}_k \cdot d\boldsymbol{\ell} = \delta_{jk}
\end{equation}
\end{definition}

\begin{definition}[VNMM 2D $\mathcal{L}^1$ Formulation]
The Vector Nodal Meshless Method (VNMM 2D) based on the solenoidal linear basis employs the 3-term polynomial space:
\begin{equation}
    \mathcal{L}^1 = \operatorname{span}\left\{ \begin{bmatrix} 1 \\ 0 \end{bmatrix}, \; \begin{bmatrix} 0 \\ 1 \end{bmatrix}, \; \begin{bmatrix} y \\ -x \end{bmatrix} \right\}
\end{equation}
Given a triad of support nodes at the edge midpoints $\{P_k = M_k, \mathbf{t}_k\}_{k=1}^3$, the VNMM shape functions $\mathbf{N}_k(\mathbf{x}) \in \mathcal{L}^1$ are uniquely defined by the pointwise Kronecker projection condition:
\begin{equation}
    \mathbf{N}_k(P_j) \cdot \mathbf{t}_j = \delta_{jk}, \quad \forall j, k \in \{1, 2, 3\}
\end{equation}
\end{definition}

\section{The Dual Equivalence Theorem}

\begin{theorem}[Dual Equivalence of Whitney 1-Forms and VNMM 2D $\mathcal{L}^1$]
Let $T \subset \mathbb{R}^2$ be an arbitrary non-degenerate triangle. Consider the support triad of vector nodes $\{P_k, \mathbf{t}_k\}_{k=1}^3$ collocated at the midpoints $P_k = M_k$ with tangent vectors $\mathbf{t}_k$. Then:
\begin{enumerate}
    \item \textbf{Constructive Isomorphism:} The VNMM 2D shape functions $\mathbf{N}_k(\mathbf{x})$ and the Whitney 1-forms $\mathbf{w}_k(\mathbf{x})$ are identical everywhere in $\mathbb{R}^2$ up to the edge length scaling factor $L_k$:
    \begin{equation}
        \mathbf{N}_k(\mathbf{x}) \equiv L_k \, \mathbf{w}_k(\mathbf{x}) \quad \Longleftrightarrow \quad \mathbf{w}_k(\mathbf{x}) \equiv \frac{1}{L_k} \mathbf{N}_k(\mathbf{x}), \quad \forall k \in \{1, 2, 3\}
    \end{equation}
    \item \textbf{Duality of Degrees of Freedom:} The integral edge circulation $\alpha_k^{\mathrm{FEM}}$ and the pointwise directional projection $e_k^{\mathrm{VNMM}} = \mathbf{E}(P_k) \cdot \mathbf{t}_k$ satisfy:
    \begin{equation}
        \alpha_k^{\mathrm{FEM}} = \oint_{e_k} \mathbf{E} \cdot d\boldsymbol{\ell} = L_k \, e_k^{\mathrm{VNMM}} = \lim_{\Delta L \to 0} \frac{L_k}{\Delta L} \int_{P_k - \frac{\Delta L}{2}\mathbf{t}_k}^{P_k + \frac{\Delta L}{2}\mathbf{t}_k} (\mathbf{E} \cdot \mathbf{t}_k) \, ds
    \end{equation}
    \item \textbf{Differential Invariants:} Both representations share identical differential properties:
    \begin{equation}
        \nabla \cdot \mathbf{w}_k \equiv 0 \quad \text{and} \quad \nabla \cdot \mathbf{N}_k \equiv 0 \quad (\text{identically solenoidal})
    \end{equation}
    \begin{equation}
        (\nabla \times \mathbf{w}_k)_z = \frac{1}{|T|} \quad \text{and} \quad (\nabla \times \mathbf{N}_k)_z = \frac{L_k}{|T|} \quad (\text{constant curl})
    \end{equation}
\end{enumerate}
\end{theorem}

\begin{proof}
\textbf{Step 1: Isomorphism of Spaces.} By direct inspection:
\begin{equation}
    \mathcal{W}(T) = \left\{ \begin{bmatrix} a_1 + \alpha y \\ a_2 - \alpha x \end{bmatrix} : a_1, a_2, \alpha \in \mathbb{R} \right\} \equiv \mathcal{L}^1
\end{equation}

\textbf{Step 2: Constancy of Tangential Component.} For any $\mathbf{u}(\mathbf{x}) \in \mathcal{W}(T) \equiv \mathcal{L}^1$, its Jacobian matrix is skew-symmetric:
\begin{equation}
    J_{\mathbf{u}} = \nabla \mathbf{u} = \begin{bmatrix} 0 & \alpha \\ -\alpha & 0 \end{bmatrix}, \quad J_{\mathbf{u}}^T = -J_{\mathbf{u}}
\end{equation}
Parameterizing edge $e_k$ by arc length $s \in [-L_k/2, L_k/2]$ centered at $M_k$: $\mathbf{x}(s) = M_k + s \mathbf{t}_k$. The directional derivative satisfies:
\begin{equation}
    \frac{d}{ds} (\mathbf{u}(\mathbf{x}(s)) \cdot \mathbf{t}_k) = \mathbf{t}_k^T J_{\mathbf{u}} \, \mathbf{t}_k \equiv 0
\end{equation}
Hence, $\mathbf{u}(\mathbf{x}) \cdot \mathbf{t}_k = \mathbf{u}(M_k) \cdot \mathbf{t}_k$ is strictly constant along edge $e_k$.

\textbf{Step 3: Reduction of Line Integral.} Integrating the constant along edge $e_j$:
\begin{equation}
    \oint_{e_j} \mathbf{w}_k \cdot d\boldsymbol{\ell} = \Big( \mathbf{w}_k(M_j) \cdot \mathbf{t}_j \Big) \int_{-L_j/2}^{L_j/2} ds = L_j \Big( \mathbf{w}_k(M_j) \cdot \mathbf{t}_j \Big)
\end{equation}
Since $\oint_{e_j} \mathbf{w}_k \cdot d\boldsymbol{\ell} = \delta_{jk}$, we have:
\begin{equation}
    L_j \Big( \mathbf{w}_k(M_j) \cdot \mathbf{t}_j \Big) = \delta_{jk} \quad \Longrightarrow \quad \Big( L_k \mathbf{w}_k(M_j) \Big) \cdot \mathbf{t}_j = \delta_{jk}
\end{equation}

\textbf{Step 4: Uniqueness.} The difference field $\mathbf{d}_k(\mathbf{x}) = \mathbf{N}_k(\mathbf{x}) - L_k \mathbf{w}_k(\mathbf{x}) \in \mathcal{L}^1$ satisfies:
\begin{equation}
    \mathbf{d}_k(M_j) \cdot \mathbf{t}_j = \delta_{jk} - \delta_{jk} = 0, \quad \forall j \in \{1, 2, 3\}
\end{equation}
Because the moment matrix $\mathbf{A}$ is non-singular for any non-degenerate triangle, the only element in $\mathcal{L}^1$ satisfying this condition is the trivial field $\mathbf{d}_k \equiv \mathbf{0}$. Therefore:
\begin{equation}
    \mathbf{N}_k(\mathbf{x}) \equiv L_k \, \mathbf{w}_k(\mathbf{x}) \quad \forall \mathbf{x} \in \mathbb{R}^2
\end{equation}

\textbf{Step 5: Differential Invariants.} The divergence vanishes identically:
\begin{equation}
    \nabla \cdot \mathbf{u} = \frac{\partial (a_1 + \alpha y)}{\partial x} + \frac{\partial (a_2 - \alpha x)}{\partial y} = 0 + 0 = 0
\end{equation}
By Stokes' theorem on $T$:
\begin{equation}
    \int_T (\nabla \times \mathbf{w}_k)_z \, d\Omega = \oint_{\partial T} \mathbf{w}_k \cdot d\boldsymbol{\ell} = \sum_{j=1}^3 \delta_{jk} = 1
\end{equation}
Since $(\nabla \times \mathbf{w}_k)_z = -2\alpha$ is constant, $(\nabla \times \mathbf{w}_k)_z = 1/|T|$, and $(\nabla \times \mathbf{N}_k)_z = L_k / |T|$.
\end{proof}

\end{document}
